\documentclass[11pt]{amsart}

\usepackage[T1]{fontenc}
\usepackage{lmodern}
\usepackage{microtype}
\usepackage{amsmath,amssymb,amsthm}
\usepackage[hidelinks]{hyperref}

\hypersetup{
  pdftitle={Cubic Graphs of Girth 11 on 138 and 142 Vertices}
}

\newtheorem{theorem}{Theorem}
\newtheorem{proposition}[theorem]{Proposition}

\newcommand{\Spec}{\operatorname{Spec}}
\newcommand{\dist}{\operatorname{dist}}

\title{Cubic Graphs of Girth 11 on 138 and 142 Vertices}
\date{}
\subjclass[2020]{05C35, 05C85}
\keywords{cubic graph, girth, order spectrum, graph excision}

\begin{document}

\begin{abstract}
We construct connected cubic graphs of girth $11$ on $138$ and $142$
vertices. The order-$142$ graph is obtained by applying the standard
two-edge subdivision construction to a recently reported order-$140$
graph. The order-$138$ graph is obtained by Biggs's excision
\cite[Theorem~4.2]{Biggs} applied to a recently reported cubic graph of
order $152$ and girth $12$; the only points requiring verification there
are that the result is connected and that its girth is $11$ rather than
$12$. Together with the
known order $140$ and the previously established spectrum from order
$144$ onward, these constructions give $N(3,11)\leq 138$. Neither
construction is new: both operations are due to earlier authors and are
applied here without modification, so the contribution is the bound.
\end{abstract}

\maketitle

A $(3,11)$-graph is a simple connected cubic graph of girth $11$.
Let $\Spec(3,11)$ denote the set of its possible orders, and let
$N(3,11)$ be the least even integer $N$ such that every even $n\geq N$
belongs to $\Spec(3,11)$. Eze, Jajcay, Jajcayov\'a, and
Z\'avack\'a~\cite{EzeEtAl} established $N(3,11)\leq 144$ and left the
even orders $128$ through $142$ unresolved. Exoo, Goedgebeur, Jooken,
Stubbe, and Van den Eede~\cite{ExooEtAl} subsequently constructed three
$(3,11)$-graphs of order $140$ and a $(3,12)$-graph of order $152$.
We apply standard insertion and excision operations to two of their
publicly available graphs. Neither operation is ours: the insertion is
Technique~1 of \cite[Section~5.1]{EzeEtAl}, and the excision is
\cite[Theorem~4.2]{Biggs}, which applies unconditionally to any cubic
graph of girth $12$, so that the order-$138$ case is a corollary of a
$1998$ theorem. What is new here is only that these two operations are
carried out on the graphs of \cite{ExooEtAl}, that the two resulting
girths are verified to be exactly $11$, and the bound
$N(3,11)\leq 138$ that follows.

\begin{theorem}\label{thm:main}
The orders $138$ and $142$ belong to $\Spec(3,11)$. Consequently,
\[
N(3,11)\leq 138.
\]
\end{theorem}

All labels below are zero-based labels in the graph6 files contained in
\path{genLifts/data/graphs/spectra.zip} in the data repository of
\cite{ExooEtAl}.

\begin{proposition}\label{prop:142}
There is a $(3,11)$-graph on $142$ vertices.
\end{proposition}

\begin{proof}
Let $G$ be the graph on the first line of
\path{spectra/n140_k3_g11.g6}. Thus $G$ is connected and cubic, has
order $140$, and has girth $11$. Put
\[
e_1=\{14,87\},\qquad e_2=\{85,99\}.
\]
Replace $e_1$ by the path $14-x-87$, replace $e_2$ by the path
$85-y-99$, and add the edge $xy$, where $x=140$ and $y=141$. Denote
the resulting graph by $G_{142}$. This is the standard two-edge
subdivision construction of \cite[Section~5.1, Technique~1]{EzeEtAl};
the girth is verified directly below rather than by appealing to the
edge-distance hypothesis stated there, and the distances used are
computed in $G-e_1-e_2$, whereas \cite{EzeEtAl} measure them in $G$.
Since
$e_1$ and $e_2$ are disjoint, $G_{142}$ is simple, connected, and cubic.

In $G-e_1-e_2$, breadth-first search gives
\[
\dist(u,v)=8
\quad\text{for all }u\in\{14,87\},\quad v\in\{85,99\}.
\]
Let $C$ be a cycle of $G_{142}$. If $C$ avoids $xy$, suppressing every
subdivision vertex used by $C$ produces a cycle of $G$. Hence
$|C|\geq 11$, and $|C|\geq 12$ whenever $C$ uses $x$ or $y$.
If $C$ uses $xy$, the remainder of $C$ contains a path in
$G-e_1-e_2$ between an endpoint of $e_1$ and an endpoint of $e_2$.
Therefore
\[
|C|\geq 1+1+8+1=11.
\]
Thus $G_{142}$ has girth at least $11$. The cycle
\[
0,3,136,127,108,70,63,114,125,138,2,0
\]
avoids $e_1$ and $e_2$ and remains an $11$-cycle in $G_{142}$.
Therefore $G_{142}$ has girth exactly $11$.
\end{proof}

\begin{proposition}\label{prop:138}
There is a $(3,11)$-graph on $138$ vertices.
\end{proposition}

\begin{proof}
Let $G$ be the graph in \path{spectra/n152_k3_g12.g6}. Thus $G$ is
connected and cubic, has order $152$, and has girth $12$. Remove
\[
S=\{0,1,2,3,136,137,144,145,146,147,148,149,150,151\}
\]
and write $H=G-S$. The vertices of degree $2$ in $H$ are precisely
\[
D=\{112,113,118,119,122,123,130,131,
     134,135,138,139,140,141,142,143\},
\]
and every other vertex of $H$ has degree $3$. Add the perfect matching
\begin{align*}
M=\{&
\{112,118\},\{113,119\},\{122,130\},\{123,131\},\\
&
\{134,140\},\{135,141\},\{138,142\},\{139,143\}
\}.
\end{align*}
Set $G_{138}=H+M$. This is precisely Biggs's excision, and not merely an
operation of that type: by \cite[Theorem~4.2]{Biggs}, if there is a
cubic graph with $n$ vertices and girth $g$, then there is a cubic graph
with $n-\epsilon(g)$ vertices and girth $g-1$, where
$\epsilon(g)=2^{r+1}-2$ for $g=4r$ or $g=4r+1$, so that
$\epsilon(12)=2^{4}-2=14$; and the excised set $S$ above is Biggs's
tree, an adjacent pair together with all vertices at distance at most
$2$ from it (cf.\ \cite[Example~4.5]{Biggs}). Applied to the order-$152$
graph of \cite{ExooEtAl}, that theorem already yields a cubic graph of
order $138$ and girth at least $11$; what the remainder of this proof
adds is only the verification that the graph so obtained is connected
and that its girth is $11$ rather than $12$. For every edge
$\{u,v\}\in M$, one has
$\dist_H(u,v)=10$. Moreover, if $u$ and $v$ are endpoints of two
distinct edges of $M$, then $\dist_H(u,v)\geq 5$. In particular, no
edge of $M$ already belongs to $H$, so $G_{138}$ is simple. The stated
degrees imply that it is cubic, and breadth-first search shows that it
is connected.

Let $C$ be a cycle in $G_{138}$ and put
\[
k=|E(C)\cap M|.
\]
If $k=0$, then $C$ is a cycle of $G$, so $|C|\geq 12$. If $k=1$,
deleting the unique matching edge from $C$ leaves an $H$-path of length
at least $10$, and hence $|C|\geq 11$. If $k\geq 2$, deleting the $k$
matching edges decomposes $C$ into $k$ paths in $H$, each joining
endpoints of distinct matching edges. Each such path has length at
least $5$, so
\[
|C|\geq k+5k=6k\geq 12.
\]
Thus $G_{138}$ has girth at least $11$. Finally,
\[
4,34,97,128,141,135,124,69,49,93,32,4
\]
is an $11$-cycle in $G_{138}$, whose only edge in $M$ is
$\{135,141\}$. Hence $G_{138}$ has girth exactly $11$.
\end{proof}

\begin{proof}[Proof of Theorem~\ref{thm:main}]
Propositions~\ref{prop:142} and~\ref{prop:138} give the two asserted
orders. Eze et al.~\cite{EzeEtAl} established that every even order at
least $144$ belongs to $\Spec(3,11)$, and Exoo et al.~\cite{ExooEtAl}
supplied order $140$. Therefore every even order at least $138$ belongs
to $\Spec(3,11)$.
\end{proof}

No assertion is made here about the orders
$128,130,132,134,$ and $136$.

\paragraph{Data and verification.}
The parent archive is available in the repository
\href{https://github.com/AGT-Kulak/smallRegGirthGraphs}%
{\texttt{AGT-Kulak/smallRegGirthGraphs}} at commit
\href{https://github.com/AGT-Kulak/smallRegGirthGraphs/commit/c2a1a4e334cd9f4dd31247f1ebfb6e18c004e4bc}%
{\texttt{c2a1a4e334cd}}. At that commit the archive
\path{genLifts/data/graphs/spectra.zip} has SHA-256 checksum
\begin{center}
\small\texttt{5fea94f6985d681adf33f4c104c391ea0f0ed8f80af82b4241500e5c84c9b3ac}
\end{center}
No verification script accompanies this note. Everything needed to
reconstruct and recheck both graphs is stated above: the two source
graphs and their locations in the archive, the edges $e_1$ and $e_2$,
the excised set $S$, the matching $M$, the relevant distances, and an
$11$-cycle in each of $G_{142}$ and $G_{138}$.

\begin{thebibliography}{9}

\bibitem{Biggs}
N.~Biggs,
\emph{Constructions for cubic graphs with large girth},
Electron. J. Combin. \textbf{5} (1998), Article A1,
\url{https://doi.org/10.37236/1386}.

\bibitem{EzeEtAl}
L.~C. Eze, R.~Jajcay, T.~Jajcayov\'a, and D.~Z\'avack\'a,
\emph{Theoretical and computational approaches to determining sets of
orders for $(k,g)$-graphs},
Discrete Appl. Math. \textbf{380} (2026), 572--587,
\url{https://doi.org/10.1016/j.dam.2025.10.049}.

\bibitem{ExooEtAl}
G.~Exoo, J.~Goedgebeur, J.~Jooken, L.~Stubbe, and T.~Van den Eede,
\emph{New small regular graphs of given girth: the cage problem and
beyond},
arXiv:2511.07247v2, 22 July 2026,
\url{https://arxiv.org/abs/2511.07247}.

\end{thebibliography}

\end{document}